% takeoff_landing.tex -- Methodology: takeoff from the floor, lift and landing (draft 4 Oct 2026).
% Code: mpc_planner/creep_controller.py (arc creep, hand-over latch), planner_node.py (pretension,
% breakaway, ff_cap_release_s, lift ramp, LAND, TouchdownDetector), bringup/config/*.yaml.
% Evidence: 23 Sep rig non-lift (start_taut), 30 Sep rig review, R0817/R0818, R0834-R0839 (B2),
% R0807-R0816 (hand-over angle in the twin), RIG-1001-* (hand-over and rods on the rig).
% Condensed 12 Oct 2026 to the plan's half page; the detail is Appendix C "Floor start and landing in detail".
% Restructured 9 Oct (structure audit): the outcomes (the hand-over angle flown, fig:flight) moved to
% Section 4.1.

\subsubsection{Takeoff, lift and landing}
\label{sec:takeoff}

The planner assumes taut rods, but every flight starts and ends with the ring on the floor and the rods slack or
flat, where the UAVs sit about 0.7\,m from the ring's centre, outside the cone the planner's references put them
on; started from rest, the planner pulled the UAVs inwards over the ring instead of lifting it, the cause of the
first failed lifts on the rig. Four steps take a flight from the floor into the regime the planner models, and three
measures take it back; Appendix~\ref{sec:app_floor} gives each with its numbers and the failure it answers.
\begin{itemize}
    \item \textbf{Rod sweep.} On TAKEOFF each UAV is anchored where it rests and its rod is swept about its own
    attachment point from its measured elevation to the planned one (60$^\circ$ since 7 October), the UAV's
    reference riding the far end of the rod, so that the rod stays straight and the UAV never crosses over the
    ring. The planner takes over once every rod has held within 8$^\circ$ of the planned elevation for 1\,s, and
    measures each rod's length from motion capture at that hand-over for its model and its tautness gate.
    \item \textbf{Pretension and breakaway.} Every UAV's share of the planned pull then rises together from zero
    over 3\,s with the ring's height target at its resting height; a step here launched the ring and released two
    magnets on the rig. When the ring rises 1\,cm or faster than 0.03\,m/s it has broken free, and the fraction of
    the planned pull reached at that moment is frozen as a cap on the trackers' feedforward, released back to the
    full pull over 2\,s once the lift starts.
    \item \textbf{Lift.} The height target ramps to the target height at 0.1\,m/s, eased in and out, and the
    trajectory starts when the ramp completes.
    \item \textbf{Landing.} On LAND the ring's height target descends at 0.2\,m/s with the formation unchanged
    until the ring rests. A rigid rod then pushes a UAV descending straight down outward as the rod lengthens
    horizontally, which tipped UAVs over in the twin, so each UAV's reference unwinds along its rod's circle about
    the grounded attachment point before it descends (twin only: the rig's tethers go slack and it lands straight
    down); the planner blends each UAV's thrust and rod feedforward to plain hover over 1\,s instead of dropping
    the rod term at touchdown; and a UAV back at its resting height with the ring down idles armed until the fleet
    disarms, so that a UAV that lands first does not track a reference below the floor. Touchdown is declared when
    the UAVs stop following the descending reference for 1\,s after at least 5\,cm of descent, and the fleet
    manager disarms the fleet.
\end{itemize}
